% ===============================================================
\chapter{Definitions or Concepts I May Have Missed}
% ===============================================================

\begin{enumerate}
  \item \textbf{Week 7.} \emph{Exclusive disjunction}: an XOR gate. Exclusive ``or''; it maps to the semantic value ``false'' if both disjuncts are true.
  \item \textbf{Week 7.} \emph{Inclusive disjunction}: inclusive ``or''; it can have the same semantic value as a conjunction.
  \item \textbf{Week 8.} \emph{Epistemological holism} is the view that beliefs are not confirmed or disconfirmed by experience individually, but rather as an entire web, body, or tribunal.
  \item \textbf{Week 7.} \emph{Donnellan--Kripke debate}: whereas Frege and Strawson held that definite descriptions are singular terms, Russell instead maintained that they are quantificational expressions. In contrast, in his paper `Reference and Definite Descriptions', Keith Donnellan argues that definite descriptions have both sorts of meanings, and so are ambiguous. However, in `Speaker's Reference and Semantic Reference', Saul Kripke argues that definite descriptions have just one, Russellian meaning, with the appearance of reference explained as a kind of implicature. His argument appeals to GR: prefer a pragmatic explanation to a semantic one.
\end{enumerate}
